\section{Introduction}\label{sec:introduction}
\PARstart{V}{isual} tyre inspection combines several different questions. A system may locate a tyre, identify its tread region, estimate image-relative boundaries, classify a usage-related category, and expose the evidence behind its output. These functions have different supervision and different meanings. A good mask is evidence of region agreement, whereas a reliable physical-wear estimate requires suitable measurements and independent evaluation. Optical tread-depth systems such as TireEye explicitly connect image structure to a physical reference \cite{tireeye}. A classifier trained on mileage categories does not acquire that connection merely by producing a plausible heat map.\AIdraftingnote

The present work grew from a compact, custom multitask system developed for an existing tyre-inspection application. Its dense outputs were useful internally, but mileage-proxy classification transferred poorly to the allocated held-out tyres. Subsequent changes addressed gradient interference, appearance sensitivity, reference exposure, spatial layout, and separation of geometry from texture learning. These revisions supplied informative failures rather than a consistently improving classifier. They also made it possible to ask a narrower question: \emph{does useful local texture evidence survive the operation that converts spatial tokens into a global representation?}

This question matters because reconstruction and representation diagnostics can disagree for structural reasons. A local reconstruction head can read individual tokens without ever using a pooled image vector. Simultaneously, an auxiliary global head can be dominated by one direction of variation. Observing both properties does not constitute a contradiction. Nor does increasing a rank statistic automatically establish better evidence: coordinate scaling changes spectral balance, and a recorded mask layout can contribute variation unrelated to the desired texture content.

We study these issues in TRACE (Texture Reconstruction with Anchored Canonical Evidence), a small model built around frozen predicted geometry and fixed local image descriptors. Its six targets describe contrast, edge energy, orientation, curvature, and contrast occupancy. They are observable image properties, not annotated groove depth or damage. The model improves local and masked reconstruction against the declared controls but never satisfies its original representation-selection requirements. We preserve that outcome and examine the saved final tokens using the same frozen decoder, rather than retraining until a favorable endpoint appears.

The contributions are threefold. First, we document a compact geometry-anchored model and the experimental context that motivated a shift from proxy classification to local texture analysis, including unsuccessful revisions and unchanged selection rules. Second, we formulate a controlled saved-feature assay that compares nine spatial aggregation interventions on identical supported targets, with training-only templates and explicit frozen-head replay accounting. Third, we connect the measured reconstruction losses to descriptor-specific effects, component-centered covariance, invertible scaling, and an exploratory mask-layout association. These analyses distinguish local predictive utility from broad global diversity and identify a constrained future design question.

The study uses 570 images from 12 user-confirmed physical tyres. It is a retrospective, limited-data investigation with only two conservative components in each held-out role. The historical test role was inspected repeatedly during development. Accordingly, its results are descriptive corroboration on a known allocation, not an unbiased estimate from a new external test. We report no new physical-wear accuracy, no independent safety decision, and no newly selected TRACE checkpoint.
